\section{Methods}

\subsection{Agent Architecture}

Each agent in HSAP is characterized by sex (assigned 50/50 at birth), age (in model time steps; 10 steps = 1 model-year), and a suite of physiological and behavioral state variables. Agents are initialized with age drawn uniformly from $[5, 40)$ and energy, health, and rank set to 1.0, 1.0, and 0.5 respectively.

\paragraph{Hormonal state.}
Each agent carries three hormonal axes: testosterone, estrogen, and cortisol. Male testosterone is initialized near 1.0 (\texttt{base\_male\_testosterone}) with Gaussian noise (SD=0.15); female testosterone near 0.3 (SD=0.05). Estrogen is initialized near 1.0 for females and 0.3 for males. Cortisol is initialized near 0.5 for both sexes (SD=0.1). All hormonal values are clamped to $[0, 1]$ at each time step.

\paragraph{Behavioral state.}
Derived from hormones and environment at each step: \texttt{aggression\_tendency}, \texttt{maternal\_defense\_tendency} (females), \texttt{mating\_drive}, \texttt{fertility}, and \texttt{offspring\_survival\_probability}. These are recomputed every time step; agents do not carry fixed behavioral types.

\paragraph{Social state.}
Each agent has a dominance rank (clamped $[0, 1]$), territory quality, \texttt{cooperation\_bonus} (accumulated from cooperative actions, capped at 0.3), and a pregnancy flag with gestation countdown (\texttt{gestation\_time} = 5 steps).

\subsection{Environment}

The environment provides four static parameters and several dynamic state variables.

\paragraph{Static parameters} (set per scenario):
\begin{itemize}
\item \texttt{resource\_abundance} (0 to 1): baseline resource availability
\item \texttt{predator\_pressure} (0 to 1): external threat level
\item \texttt{disease\_pressure} (0 to 1): disease prevalence
\item \texttt{space\_constraint} (0 to 1): controls how quickly territory availability declines with density
\item \texttt{carrying\_capacity} (integer): $K$ for density computation
\end{itemize}

\paragraph{Dynamic state} (updated each step):
\begin{itemize}
\item \texttt{resource\_abundance}: regenerates toward maximum, consumed by population; modulated by seasonality
\item \texttt{territory\_availability}: $\max(0, 1 - \rho \cdot S)$, where $\rho = N / K$
\item \texttt{predation\_risk}: $P \cdot (1 + 0.5 \cdot \rho)$
\item \texttt{disease\_risk}: $D \cdot (1 + 0.5 \cdot \rho)$
\end{itemize}

\subsection{Endocrine Update}

Hormones are recomputed for every agent at each time step, creating the causal feedback loop that defines HSAP.

\paragraph{Male testosterone} is driven by five factors:

\begin{equation}
T_{m,t+1} = \min\left(T_{\max},\; T_{\text{base}} + 0.3 \cdot \text{rank} + 0.4 \cdot M_d - 0.3 \cdot \rho - 0.4 \cdot (1 - P) - 0.2 \cdot (1 - \text{health})\right)
\label{eq:male_T}
\end{equation}

The term $-0.4 \cdot (1 - P)$ is the critical endocrine feedback: when predator pressure is high, the term is near zero (testosterone is not reduced); when predator pressure is low, the term subtracts up to 0.4 (testosterone is downshifted). This is the mechanism through which threat reduction propagates to behavioral change.

\paragraph{Female testosterone} responds to resource scarcity and density:

\begin{equation}
T_{f,t+1} = \min\left(T_{\max},\; 0.3 + 0.15 \cdot \text{rank} + 0.4 \cdot (1 - R) + 0.3 \cdot \rho\right)
\label{eq:female_T}
\end{equation}

\paragraph{Cortisol} (both sexes) integrates density, rank, health, predation, and resource scarcity:

\begin{equation}
C = 0.5 + 0.1 \cdot \rho - 0.3 \cdot \text{rank} + 0.5 \cdot (1 - \text{health}) + 0.2 \cdot P + 0.3 \cdot (1 - R)
\label{eq:cortisol}
\end{equation}

\paragraph{Estrogen} (females) is suppressed by cortisol:

\begin{equation}
E = 1.0 - 0.3 \cdot C
\label{eq:estrogen}
\end{equation}

\subsection{Behavioral Output}

From hormones, agents derive:

\begin{itemize}
\item \textbf{Aggression}: base $= T \cdot 0.5 + C \cdot 0.2$, with sex-specific additions (females: offspring defense + resource competition; males: dominance bonus)
\item \textbf{Fertility}: base $= 0.5$, modified by hormones, age, density-dependent reproductive restraint (activates above density 0.7), and territory availability
\item \textbf{Mating drive}: males $= T \cdot 0.5 + 0.3$; females $= E \cdot 0.4 + 0.2$
\item \textbf{Offspring survival probability}: affected by testosterone, cortisol, density, and resource scarcity
\end{itemize}

Agents select actions stochastically from: forage, dominate, disperse, withdraw, cooperate, guard\_offspring. Action probabilities depend on current behavioral state and environmental conditions.

\subsection{Reproduction}

Mating occurs when a randomly selected male (probability proportional to mating\_drive) encounters a fertile female above a fertility gate (default 0.3). Both must have energy above a threshold (0.4) and be in reproductive age (10--60 model-years). On successful mating, the female becomes pregnant with gestation time of 5 steps.

At birth, litter size is drawn from a normal distribution (mean=4.0, SD=1.0, floored to 1). Each pup survives with probability determined by the mother's \texttt{offspring\_survival\_probability}. The mother pays an energy cost of 0.3. Newborns face an additional mortality check.

\paragraph{Infanticide} triggers when density exceeds 0.8, with probability proportional to agent aggression and density above threshold. \paragraph{Offspring neglect} triggers when density exceeds 0.7 or cortisol exceeds 0.7, removing one offspring.

\subsection{Mortality}

Death probability per time step is computed as:

\begin{equation}
P_{\text{death}} = \min\left(1,\; (P_{\text{base}} + P_{\text{pred}} + P_{\text{disease}} + P_{\text{starvation}} + P_{\text{injury}} + P_{\text{old}} + P_{\text{crowd}}) \times \tau \times M\right)
\label{eq:mortality}
\end{equation}

where $\tau = 0.1$ (age increment per step) and $M$ is the mortality multiplier (1.0 normally, 0.4 during post-sink recovery).

\begin{table}[h]
\centering
\caption{Mortality components}
\label{tab:mortality}
\begin{tabular}{ll}
\toprule
Component & Formula \\
\midrule
Base & 0.02 (constant) \\
Predation & $0.05 \cdot P_{\text{risk}} \cdot (1 - 0.02 \cdot (1-R)) \cdot \text{modifiers}$ \\
Disease & $0.03 \cdot D_{\text{risk}} \cdot (1 + 0.5 \cdot (1 - \text{health})) \cdot (1 + 0.5 \cdot \rho)$ \\
Starvation & If energy $< 0.2$: $0.3 \cdot (1-R) \cdot (1 - \text{energy}/0.2)$ \\
Injury & 0.1 if injured \\
Old age & $0.01 \cdot (\text{age} - 50)$ if age $> 50$ \\
Crowding & $0.05 \cdot \max(0, \rho - 0.8) + \eta \cdot M_{\text{sink}}$ \\
\bottomrule
\end{tabular}
\end{table}

\subsection{Behavioral Sink and Recovery}

The behavioral sink is a density-driven phase transition that suppresses reproductive output even after density falls below carrying capacity.

\paragraph{Activation.}
When density exceeds \texttt{behavioral\_sink\_on\_threshold} (default 0.75), the sink engages. A \texttt{sink\_factor} is computed as a continuous measure of sink severity:

\begin{equation}
\eta = \min\left(1,\; \frac{\rho - 0.75}{1 - 0.75}\right)
\label{eq:sink_factor}
\end{equation}

\paragraph{Effects while active.}
Fertility is multiplicatively suppressed: $F \times (1 - \eta \cdot 0.85)$. Mating drive is suppressed: $M_d \times (1 - \eta \cdot 0.4)$. Withdrawal behavior and offspring neglect probabilities are increased.

\paragraph{Deactivation (hysteresis).}
The sink disengages only when density falls below \texttt{behavioral\_sink\_off\_threshold} (default 0.50), creating a hysteresis band (on at 0.75, off at 0.50). A minimum duration of 30 steps must elapse before deactivation is permitted. An optional auto-recovery timer forces disengagement after a fixed number of steps.

\paragraph{Post-sink recovery.}
After sink disengagement, a recovery phase lasting 100 steps begins. During recovery: fertility and mating drive are multiplicatively boosted by 1.3; mortality is reduced to 40\% of normal. If population drops below 30 during recovery, refugees (up to 10 per step) are injected to supplement the depleted population.

\subsection{Metrics}

Per time step, HSAP records: population size, density, sex ratio, resource abundance, predator pressure, disease pressure, territory availability, age structure (juvenile/adult/senior counts), mean hormones (testosterone, estrogen, cortisol, male\_T, female\_T), behavioral outputs (male/female aggression, female defense), reproductive metrics (fertility, births, deaths, matings, pregnancies, infanticide, neglect, refugees), sink state (active, post-recovery, sink\_factor), external threat index, and the composite HSAP index.

The HSAP index is a weighted composite of five signals: low external threat (25\%), low male aggression (20\%), high female aggression (20\%), low fertility (20\%), and population stability (15\%), with a viability dampener for populations below 50.

\paragraph{Summary metrics} (computed over the full run): final population, peak population, minimum population, crash ratio (final/peak), time to stability, mean male/female testosterone, mean male/female aggression, mean fertility, mean cortisol, extinction status, and seed number.
